\itemhead{p.103 练习 · 第1题（教材p.103）}{XJ-B1-C04-PR-P103-Y0797}
\label{XJ-B1-C04-PR-P103-Y0797}
{\footnotesize 湘教版教材p.103 / PDF第110页 / p.103 练习 / 第1题}
\field{原题}{结合图4.1-1中幂函数 $y=x$，$y=x^2$，$y=x^3$，$y=\sqrt x$，$y=\dfrac1x$ 的图象，回答下列问题：
\subq{1}这些函数中，哪些是奇函数，哪些是偶函数？
\subq{2}这些函数的定义域、值域分别是什么？
\subq{3}这些函数的单调区间分别是什么？
\subq{4}通过上述问题，你有什么发现？
\bankfig{C04-PR103}{0.63}}
\field{参考答案}{（1）奇函数：$x,x^3,1/x$；偶函数：$x^2$；$\sqrt x$ 非奇非偶。（2）$x,x^3$ 的定义域值域均为 $\mathbb R$；$x^2$ 为 $\mathbb R,[0,+\infty)$；$\sqrt x$ 均为 $[0,+\infty)$；$1/x$ 均为 $\mathbb R\setminus\{0\}$。（3）$x,x^3$ 全域递增；$x^2$ 在 $(-\infty,0]$ 递减、$[0,+\infty)$ 递增；$\sqrt x$ 在其定义域递增；$1/x$ 在两个半轴分别递减。（4）正半轴上正指数幂递增、负指数幂递减；全域性质还要看具体指数和定义域。}
{\small \field{关联知识}{P03 幂函数的概念、图象与性质（教材帮p.125、126、127 / PDF134、135、136）}
\field{方法与易错点}{设 $f(x)=x^\alpha$，用已知点求指数，再查该指数的定义域、奇偶性和各单调区间。 易错：把 $3x^2$ 当标准幂函数；将第一象限性质直接用于负半轴；分数负指数遗漏0禁点。}
\field{答案依据}{答案册相关题组：PDF第30、31页（册内页29、30）；按当前题库原题逐项重算，不按旧版题号机械对套。}}
